\documentclass[../../main.tex]{subfiles}
\begin{document}

\begin{flushright}
\footnotesize\textit{【自動文字起こし・要確認】}
\end{flushright}

[解] $k \in \mathbb{N}$ に対し、$a_k = \int_{\dfrac{k-1}{n}\pi}^{\dfrac{k}{n}\pi} e^{-x} |\sin nx| dx$ とおく。$A_n = \int_0^\pi e^{-x} |\sin nx| dx$
おくと、
\[ A_n = \sum_{k=1}^n a_k \cdots \text{\textcircled{1}} \]
である。
\begin{eqnarray*}
a_{k+1} &=& \int_{\dfrac{k}{n}\pi}^{\dfrac{k+1}{n}\pi} e^{-x} |\sin nx| dx \\
&=& \int_{\dfrac{k-1}{n}\pi}^{\dfrac{k}{n}\pi} e^{-(t+\dfrac{\pi}{n})} |\sin n(t+\dfrac{\pi}{n})| dt\ (t = x - \dfrac{\pi}{n}) \\
&=& e^{-\dfrac{\pi}{n}} a_k \cdots \text{\textcircled{2}}
\end{eqnarray*}
又、
\begin{eqnarray*}
a_1 &=& \int_0^{\dfrac{\pi}{n}} e^{-x} \sin nx dx \\
&=& \left[ \dfrac{e^{-x}}{1+n^2} (-\sin nx - n \cos nx) \right]_0^{\dfrac{\pi}{n}} \\
&=& \dfrac{1}{1+n^2} \{ e^{-\dfrac{\pi}{n}} (n) - 1(-n) \} \\
&=& \dfrac{n}{1+n^2} (e^{-\dfrac{\pi}{n}} + 1) \cdots \text{\textcircled{3}}
\end{eqnarray*}
だから \textcircled{2} をくり返し用いて、
\[ a_k = e^{-\dfrac{(k-1)}{n}\pi} a_1 \]
\textcircled{1} に代入して
\begin{eqnarray*}
A_n &=& \sum_{k=0}^{n-1} a_1 e^{-\dfrac{k}{n}\pi} \\
&=& a_1 \dfrac{1-e^{-n \dfrac{\pi}{n}}}{1-e^{-\dfrac{\pi}{n}}} = \dfrac{n}{1+n^2} (1+e^{-\dfrac{\pi}{n}}) \dfrac{1-e^{-\pi}}{1-e^{-\dfrac{\pi}{n}}} \\
&=& \dfrac{n}{1+n^2} \dfrac{\dfrac{\pi}{n}}{1-e^{-\dfrac{\pi}{n}}} \left(\dfrac{1}{\dfrac{\pi}{n}}\right) (1+e^{-\dfrac{\pi}{n}})(1-e^{-\pi}) \\
&=& \dfrac{1}{1+\dfrac{1}{n^2}} \dfrac{\dfrac{\pi}{n}}{1-e^{-\dfrac{\pi}{n}}} (1-e^{-\pi}) \left(\dfrac{1}{\pi}\right) (1+e^{-\dfrac{\pi}{n}}) \\
&\xrightarrow{n\to\infty}& 1 \cdot 1 \cdot (1-e^{-\pi}) \cdot \left(\dfrac{1}{\pi}\right) \cdot 2 = \dfrac{2}{\pi}(1-e^{-\pi})
\end{eqnarray*}

[解2] (\textcircled{1} まで同じ)
$[\dfrac{k-1}{n}\pi, \dfrac{k}{n}\pi]$ において、$e^{-\dfrac{k}{n}\pi} \le e^{-x} \le e^{-\dfrac{k-1}{n}\pi}$ だから、$|\sin nx| \ge 0$ とあわせて、
\[ e^{-\dfrac{k}{n}\pi} \int_{\dfrac{k-1}{n}\pi}^{\dfrac{k}{n}\pi} |\sin nx| dx \le a_k \le e^{-\dfrac{k-1}{n}\pi} \int_{\dfrac{k-1}{n}\pi}^{\dfrac{k}{n}\pi} |\sin nx| dx \]
\[ \dfrac{2}{n} e^{-\dfrac{k}{n}\pi} \le a_k \le \dfrac{2}{n} e^{-\dfrac{k-1}{n}\pi} = e^{\dfrac{\pi}{n}} \left( \dfrac{2}{n} e^{-\dfrac{k}{n}\pi} \right) \cdots \text{\textcircled{2}} \]
ここで、$B_n = \sum_{k=1}^n \dfrac{2}{n} e^{-\dfrac{k}{n}\pi}$ とすると、
\[ B_n = \dfrac{2}{n} e^{-\dfrac{\pi}{n}} \dfrac{1-e^{-\pi}}{1-e^{-\dfrac{\pi}{n}}} \]
だから \textcircled{1}、\textcircled{2} より、
\[ B_n \le A_n \le e^{\dfrac{\pi}{n}} B_n \]
はさみうちから
\[ A_n \to B_n \]
で、
\[ B_n = \dfrac{\dfrac{\pi}{n}}{1-e^{-\dfrac{\pi}{n}}} \dfrac{2}{\pi} e^{-\dfrac{\pi}{n}} (1-e^{-\pi}) \to \dfrac{2}{\pi}(1-e^{-\pi}) \]
から
\[ A_n \to \dfrac{2}{\pi}(1-e^{-\pi}) \]

\end{document}
